% !TeX root = ../main.tex
\section{Results}
\label{sec:results}

We evaluate all methods on the held-out FHM target set under the
leakage-controlled protocol in Section~\ref{sec:setup:protocol}.
Trainable models with a complete multi-seed result set are reported as
mean~$\pm$~standard deviation, whereas deterministic vision--language
baselines and canonical-seed ablations are reported as single runs.
Structured results use schema-validated FHM silver annotations and are
therefore interpreted separately from the original harmfulness labels.

We organize the empirical study around five questions:
\textbf{RQ1} whether the proposed framework improves cross-domain
harmfulness detection;
\textbf{RQ2} whether structured interpretations transfer from HarMeme to
FHM;
\textbf{RQ3} whether external knowledge is useful and whether verification
adds value beyond retrieval alone;
\textbf{RQ4} which acquired candidates are retained or rejected by the
verifier; and
\textbf{RQ5} where the framework still fails under domain shift.
RQ1--RQ2 are addressed in this section, while RQ3--RQ5 are analyzed in
Section~\ref{sec:analysis}.

\subsection{Overall Cross-Domain Performance}
\label{sec:results:harmfulness}

Table~\ref{tab:main-results} compares harmfulness detection after
development exclusively on the HarMeme source domains.
The comparison includes unimodal and simple multimodal controls,
specialized harmful-meme models, general-purpose vision--language
baselines, and our evidence-grounded framework.

\begin{table*}[t]
    \caption{
    Held-out FHM harmfulness detection under the locked
    HarMeme-to-FHM protocol.
    Multi-seed trainable methods are reported as mean~$\pm$~standard
    deviation; deterministic inference-only conditions are reported as
    single runs.
    }
    \label{tab:main-results}
    \centering
    \renewcommand{\arraystretch}{1.18}
    \begin{tabular*}{0.88\textwidth}{@{\extracolsep{\fill}}lcc@{}}
\toprule
Method & Macro-F1 & AUROC \\
\midrule
\multicolumn{3}{l}{\textit{Unimodal / simple multimodal}} \\
Text-only DeBERTa      & $0.388 \pm 0.000$ & $0.000 \pm 0.000$ \\
Image-only OpenCLIP    & $0.388 \pm 0.000$ & $0.499 \pm 0.012$ \\
Image--Text Concat     & $0.388 \pm 0.000$ & $0.000 \pm 0.000$ \\
OpenCLIP Classifier    & $0.396 \pm 0.018$ & $0.502 \pm 0.020$ \\
\addlinespace
\multicolumn{3}{l}{\textit{Specialized harmful-meme models}} \\
HateCLIPper             & $0.600 \pm 0.004$ & $0.653 \pm 0.003$ \\
MemeCLIP                 & $0.580 \pm 0.014$ & $0.643 \pm 0.009$ \\
RGCL                     & $0.564 \pm 0.021$ & $0.609 \pm 0.012$ \\
RA-HMD                   & $0.648 \pm 0.005$ & $0.691 \pm 0.009$ \\
\addlinespace
\multicolumn{3}{l}{\textit{General-purpose VLM}} \\
Qwen2.5-VL Direct        & $0.444$ & $0.745$ \\
Qwen2.5-VL Rationale     & $0.627$ & $0.734$ \\
% Pending: Qwen2.5-VL RAG is still running. This row must not survive the
% final result freeze with missing values.
Qwen2.5-VL RAG           & -- & -- \\
\addlinespace
\multicolumn{3}{l}{\textit{Proposed}} \\
% Pending: Ours requires the complete audited five-seed result set. This row
% must not survive the final result freeze with missing values.
Ours                     & -- & -- \\
\bottomrule
\end{tabular*}

\end{table*}

Among the completed trainable comparisons, RA-HMD obtains the highest mean Macro-F1
($0.648\pm0.005$) and the highest mean AUROC ($0.691\pm0.009$).
All specialized-model values are our common-protocol reproductions using
official-code adapters, not scores copied from the original papers.

The deterministic Qwen2.5-VL conditions show a different metric profile.
\textsc{Direct} obtains 0.444 Macro-F1 and 0.745 AUROC, while
\textsc{Rationale} obtains 0.627 Macro-F1 and 0.734 AUROC.
Thus, rationale prompting substantially raises hard-label Macro-F1 in this
locked comparison, whereas both conditions exceed the completed trainable
models in AUROC.  The \textsc{RAG} condition and the complete five-seed
results for our framework are still pending, so we defer overall ranking and
final interpretation until the comparison table is complete.

\subsection{Structured Interpretation}
\label{sec:results:structured}

RQ2 asks whether the model transfers more than a binary harmfulness
decision.
We therefore evaluate target presence, target granularity, primary intent,
rhetorical tactic, and multimodal relation on FHM samples with valid silver
supervision.
The five structured fields capture complementary aspects of meme
interpretation: whether a target is present, the granularity of that target,
the primary communicative intent, the rhetorical strategies expressed by the
meme, and the semantic relation between image and text.
Because these annotations are schema-validated silver labels rather than
independently adjudicated gold annotations, we report each field separately
and evaluate only samples marked valid for that field.
This prevents missing or uncertain structured supervision from being treated
as an incorrect prediction.

\begin{table*}[t]
    \caption{
    Structured interpretation on schema-validated FHM silver annotations.
    Metrics are computed only over samples with valid field-level
    supervision; valid-sample counts and coverage are reported in the
    supplementary material.
    }
    \label{tab:structured-results}
    \centering
    \begin{tabular*}{0.95\textwidth}{@{\extracolsep{\fill}}lccccc@{}}
\toprule
Method & Target-P & Target-G & Intent & Tactic & MM-Rel\\
\midrule
Not run & -- & -- & -- & -- & -- \\
\bottomrule
\end{tabular*}

\end{table*}

% TODO_RESULT:
% After result freeze, summarize:
% - which structured fields transfer most reliably;
% - which fields are hardest under domain shift;
% - whether the task-specific methods and VLM baselines show different
%   strengths across fields;
% - coverage caveats where relevant.
%
% Main table target schema:
% Method | Target-P | Target-G | Intent | Tactic | MM-Rel
% Use Macro-F1 in the main paper; move tactic Micro-F1, exact match, and
% per-label results to the supplementary material.
